\section{Introduction}\label{sec:intro}

\subsection{The problem}\label{sec:intro-problem}

A $d\times d$ matrix with entries of $b(d)=\max(1,\lceil\log_2d\rceil)$ bits is written on the input tape of a
deterministic multitape Turing machine, row after row, as a word of $L=d^2b(d)$ bits. The machine has to write the
transpose. Statement~(T) says that some machine does this within $c(L+1)$ steps for a constant~$c$. The best known
machines need of the order of $L\log d$ steps, apart from a small saving in the exponent of the logarithm claimed in an
unrefereed preprint~\cite{OpenAI2026mult}, and no lower bound above linear is known for any explicit problem on machines with two work tapes. A
linear-time multiplier of integers would give a linear-time transposer~\cite{HvdH2025}, so a proof that (T) fails would
also bound multiplication from below. The formal statement is the Lean proposition
\lean{Transpose.LinearTimeTransposable} of the first paper~\cite{FirstPaper2026}, which is frozen: this paper proves
theorems about it and does not change it.

\subsection{Two conjectures}\label{sec:intro-two}

The first paper replaced (T) by a statement about codes. The run of a machine is drawn as a time--space diagram: the
positions of the heads at every step. Cutting the tapes into blocks gives a network whose nodes are the visits of a
head to a block, and a linear-time transposer yields a network code on it that delivers a piece of every row block to
every column block. The \emph{Tape Coding Conjecture} says that no such code exists on the block network of any
diagram of linear length; it implies that (T) fails~\cite[Conjecture~9.1 and Theorem~9.2]{FirstPaper2026}. Its
statement has no transition table in it, but it keeps the diagram: time, the shape of tapes, and the capacity of
every link.

The conjecture of this paper keeps none of these. Let $\rho$ assign to every column index a set of at most $w$ row
indices. A set $C$ of $d\times d$ arrays is a \emph{whole-row code for~$\rho$} if two arrays of $C$ that agree on the
rows $\rho(j)$ agree on column~$j$, for every~$j$ (\cref{def:wholerow}). For $\rho(j)=\{j\}$ the symmetric arrays are
such a code: column $j$ is row~$j$. They have $d(d+1)/2$ free entries, a little more than half. We call the fraction
of free entries the \emph{rate}.

\begin{conjecture}[Matrix Coding Conjecture]\label{conj:MCC}
There is a $\delta>0$ such that for every $w\ge1$ there is a $d_0$ with the following property. For every $d\ge d_0$
and every $\rho$ of window at most~$w$, every whole-row code for~$\rho$ with entries of $b(d)$ bits has at most
$2^{(1-\delta)L}$ elements.
\end{conjecture}

The conjecture asks only that the rate stay away from one. We expect much more, namely that the rate tends to one half
at every fixed window (the \emph{half-rate conjecture}, \cref{conj:half}), which is the form in which the synthesis
paper of the series asked the question for a fixed alphabet~\cite[Problem~9.4(c)]{SeventeenthPaper2026}. What is known
is the bound $w/(w+1)$ for the rate, and one half for $w\le2$~\cite{SeventeenthPaper2026}.

\subsection{Results}\label{sec:intro-results}

The first three theorems are the reduction. The other four are about the conjecture.

\begin{itemize}[leftmargin=1.5em]
\item \textbf{Theorem~W (whole rows; \cref{thm:W}).} In any run of linear length, a cut of bounded depth can be
extended, by revealing at most $\eps L$ further bits, so that the input bits on which the outputs of all but a small
fraction of the rows of the transpose depend lie in boundedly many rows of the input, the bound depending on $\eps$
and on the resources of the machine and not on~$d$. The theorem is about the motion of the heads alone: it assumes
neither that the machine is oblivious nor that it is correct. For oblivious machines, whose runs have cheap cuts of
bounded depth, it follows that the Matrix Coding Conjecture excludes linear-time transposers (\cref{cor:W-obl}). This is what turns a statement about runs into a statement about arrays.
\item \textbf{Theorem~T (tape coding; \cref{thm:T}).} The Matrix Coding Conjecture implies the robust form of the Tape
Coding Conjecture. With the reduction of the first paper it follows that the Matrix Coding Conjecture implies that
(T) fails, for all machines.
\item \textbf{Theorem~F (the weak form; \cref{thm:F}).} A weak form of the conjecture (\cref{conj:weak}) already
implies that (T) fails. It asks for a bound at one size $d$ beyond every threshold, and lets the deficit decrease
with the window; \cref{rem:weak-demand} says how fast it may decrease.
\item \textbf{Theorem~N (normal form; \cref{thm:N}).} Every whole-row code $C$ contains, after a translation, a code
with at least $\card C^2/2^{L}$ elements in which rows are functions of columns as well, and the relation may be taken
symmetric and reflexive at a small cost. So the conjecture may be attacked for such codes only.
\item \textbf{Theorem~O (order; \cref{thm:O}).} The largest acyclic set of entries has exactly $d^2$ minus the least
cost of an order of the rows, and every acyclic set lies in the one given by some order. Recovery rules of
transported copies of a code do not improve this bound, and on relations with a mixing property it leaves a deficit
that tends to zero as the window grows (such relations exist by \cref{prop:mixing}, which is not formalized). So no argument of this kind proves the conjecture.
\item \textbf{Theorem~E (entropy; \cref{thm:E}).} On the incidence relation of the Fano plane every polymatroid that
obeys the recovery constraints has value at most $119/3$, a code attains this, and the best acyclic bound is~$41$. A
$9\times6$ grid relation without mutual triangles has value at most $204/5$ against~$41$. Every lift of the Fano
relation whose helper rows respect the incidences and which mixes at one percent has a gap of more than $d^2/100$
between the two bounds. Such lifts exist with both degrees at most $3600$ and girth tending to infinity
(\cref{prop:lifts}, not formalized). So bounds beyond order exist at a fixed window; these relations already have a
constant deficit from order.
\item \textbf{Theorem~I (invariance; \cref{thm:I}).} A whole-row code that is invariant under the shifts of all but
$b$ pairs of a pairing of the entries has at most $q^{d^2/2+b+o(d^2)}$ elements over an alphabet of $q$ symbols, with
an explicit threshold that does not depend on the alphabet or on the relation. So a code of rate above one half must
break the symmetry of a pairing on a positive fraction of the entries.
\end{itemize}

\subsection{How the paper builds on the series}\label{sec:intro-series}

\paragraph{The first paper} supplies the frozen statement, the Tape Coding Conjecture, and the reduction from machines
that are correct on a fraction $2^{-O(d\log d)}$ of the inputs to codes on block networks, which is what makes the
robust form of the conjecture reach machines that are not oblivious~\cite{FirstPaper2026}. \Cref{thm:T} is stated
against that conjecture as it stands there.

\paragraph{The fourth and fourteenth papers} supply the cut. The fourth paper showed that for oblivious machines two
statements together exclude a transposer: every run can be cut cheaply to bounded depth, and index codes on the
resulting systems of cones are long~\cite[Theorem~4.10]{FourthPaper2026}. The fourteenth paper proved the first of the
two for all runs~\cite[Theorem~7.5]{BoundedStack2026}. \Cref{thm:W} starts from that cut, and \cref{thm:T} applies it to a block
network after lifting the network to a walk.

\paragraph{The seventeenth paper,} the synthesis, reduced the oblivious case to one function, the best rate of a
storage code on the cone relations left by cheap cuts, and described the statement that would decide it
(\cref{sec:intro-diagnosis}). It also introduced the codes of this paper, under the name of $\rho$-codes, proved the
bounds $w/(w+1)$ and one half for $w\le2$, and asked for one half at every window~\cite{SeventeenthPaper2026}. The
present paper shows that these codes, and a far weaker bound than one half, are all that the question about machines
needs.

\paragraph{Other papers.} The storage-code bounds of the fifth and seventh papers are used for the classes in which
the conjecture is known and for the certificates of \cref{thm:E}~\cite{DirectedRecovery2026,SeventhPaper2026}. The
thirteenth and fifteenth papers constructed storage codes of rate tending to one on abstract relations of bounded
degree~\cite{StorageCodes2026,HighRate2026}; whether such codes can be whole-row codes is one of the open questions of
\cref{sec:open}. The eighteenth paper treated machines that are not oblivious through their own
runs~\cite{AdaptiveTransfer2026}; \cref{thm:T} reaches them through the first paper instead.

\subsection{The diagnosis of the synthesis, point by point}\label{sec:intro-diagnosis}

The synthesis listed five properties that a statement deciding the oblivious case must
have~\cite[Section~4.4]{SeventeenthPaper2026}. We go through them for \cref{conj:MCC}.
\begin{enumerate}[label=(\arabic*),leftmargin=2em]
\item \emph{It must tie what is known to where it can be, for all schedules.} After \cref{thm:W} the place is a set of
whole rows, and the schedule has disappeared from the statement: the conjecture quantifies over all~$\rho$.
\item \emph{It must fail for abstract separated relations, on which codes of rate tending to one exist.} The
conjecture says nothing about such relations. It concerns whole-row relations, in which an entry is helped by $wd$
others, all of them in $w$ rows. No code of rate above one half is known on these at a fixed window as $d$ grows.
\item \emph{It must fail for motion that ignores the data.} The conjecture is not about motion. The codes of
\cref{thm:W,thm:T} are extracted from correct outputs, and a machine that moves its heads and transposes nothing yields
none.
\item \emph{It must constrain all scales at once.} The conjecture is a statement at one scale, and the scales enter
through the window: a cheaper cut leaves deeper cones and a larger window, so the deficit must not vanish as the
window grows. \Cref{thm:F} makes the exchange rate explicit.
\item \emph{It must use the shape of tapes,} because a bound on coding gains for graphs of bounded degree in general
would be a circuit lower bound. The shape of tapes is used in the reduction: the cheap cut exists because heads walk,
and cones meet few rows because the input is read row by row. What is left in the statement is the shape of the
array: a row and a column meet in one entry.
\end{enumerate}
So the diagnosis is met in a different way than the synthesis expected in points~(4) and~(5), and the restriction to
oblivious machines, under which it was made, is removed. The wall has not fallen. It has been moved to a statement
about arrays, and \cref{thm:O,thm:E,thm:I} say what is known on the other side.

\subsection{Questions of earlier papers}\label{sec:intro-questions}

\begin{itemize}[leftmargin=1.5em]
\item The fourth paper remarked that no implication was known, in either direction, between the plain Tape Coding
Conjecture and the pair of statements of its decomposition~\cite[Remark~4.12]{FourthPaper2026}. That question stays
open. What is new is one statement that implies both: \cref{conj:MCC} gives the bound on index codes at every depth
(\cref{cor:W-obl}) and the Tape Coding Conjecture (\cref{thm:T}).
\item The synthesis asked whether whole-row codes have rate one half at every window, over a fixed
alphabet~\cite[Problem~9.4(c)]{SeventeenthPaper2026}. We do not answer this. \Cref{thm:T,thm:F} say what any bound
below one, at the alphabet of the transposition problem, would give, and \cref{thm:I} proves one half for codes with
the symmetry of a pairing.
\item The synthesis observed that for small windows the acyclic-subgraph bound might still suffice, and that
computations indicated that for large windows it does not~\cite[Section~9]{SeventeenthPaper2026}. \Cref{thm:O} proves
the second half: on mixing relations the best acyclic bound has rate tending to one as the window grows. Window three
is still open.
\end{itemize}

\subsection{What is left}\label{sec:intro-left}

One statement: \cref{conj:weak}, or the stronger \cref{conj:MCC}, or the still stronger \cref{conj:half}. The weak form
lets the deficit decrease as the window grows, but only very slowly (\cref{rem:weak-demand}), while the bounds that
hold for all codes give a deficit of $1/(w+1)$. The smallest open case
of the half-rate conjecture is window three on relations whose rows and columns are spread. No converse is known: a
family of whole-row codes of rate tending to one would refute the conjecture and would not be a machine.
\Cref{sec:construction} says what such a family, and what a machine, would have to look like, and \cref{sec:open}
states the open problems.

\subsection{Formal status, and how to read the paper}\label{sec:intro-formal}

\Cref{thm:W,thm:T,thm:F,thm:N,thm:O,thm:I} and the finite and conditional parts of \cref{thm:E} are formalized in
Lean~4 with Mathlib~\cite{Mathlib2026}; each carries a line naming its declarations, and \cref{sec:leanindex} lists
them. They use the axioms \lean{propext}, \lean{Classical.choice} and \lean{Quot.sound} only. \Cref{thm:T,thm:F}
conclude the frozen proposition of the first paper. Not formalized, and proved in the text: the two existence statements
\cref{prop:mixing,prop:lifts}, the exclusion of \cref{prop:hyperfinite}, and the code that attains the Fano optimum.
Remarks are not formalized either: in particular the arithmetic of \cref{rem:weak-demand} and the sketch of
\cref{rem:markov}; \cref{prop:known} and the results in the table of \cref{sec:invariance} are cited from earlier
papers with their status there. A
reader who wants only the reduction can read \cref{sec:prelim,sec:rows,sec:tape,sec:weak}; a reader who wants to work
on the conjecture can start at \cref{sec:normal}.

\subsection{Related work}\label{sec:intro-related}

\paragraph{Codes.} Whole-row codes are storage codes, or equivalently index codes with side information, on a
particular family of relations; for these notions see~\cite{BBJK2011,Riis2007}. For storage codes on graphs it is
known that the simplest bounds are not always attained and that natural conjectures about optimal codes
fail~\cite{ChristofidesMarkstrom2011,Baber2017}; \cref{thm:E} is a statement of this kind for whole-row relations. In
coding theory a \emph{matrix code} usually means a code in the rank metric and an \emph{array code} a code for
arrays of disks. The conjecture of this paper is named after the matrix that is transposed and concerns neither.

\paragraph{Cuts.} Fixing the values on a cut of a computation of small depth goes back to Valiant's depth
reduction~\cite{Valiant1977}; Paul, Pippenger, Szemer\'edi and Trotter cut the computations of multitape
machines~\cite{PPST1983}. With their technique Mittal and Raz showed that a function with a static property, block
rigidity, cannot be computed in linear time by a multitape machine, even with advice, and that an explicit such
function would follow from a strong parallel repetition theorem~\cite{MittalRaz2021}; a conjecture of theirs about a
game that predicts a transpose is known to fail~\cite{Mittal2025}. Our reduction has the same shape, a static
statement that excludes linear time through a cut, with a different statement: it is about codes on arrays, and it
is specific to transposition.

\paragraph{Network coding.} Adler, Harvey, Jain, Kleinberg and Lehman related transposition in the input--output
model to network coding~\cite{AHJKL2006}, and the network coding conjecture of Li and Li~\cite{LiLi2004} would have
implied lower bounds for sorting in external memory~\cite{FHLS2019}, for multiplication
circuits~\cite{AFKL2019} and for oblivious machines with two tapes~\cite{AHJKL2006}. That conjecture has been
refuted~\cite{BravermanHe2026,ZhangLiLi2026}. The Tape Coding Conjecture of the first paper restricts it to the
networks of tapes; \cref{conj:MCC} is not a statement about networks at all. Lower bounds for transposition are known
on machines with one work tape~\cite{DMS1991,DietzfelbingerMaass1993}.

\paragraph{What we did not find.} We searched for the whole-row relations, under names such as array codes with
recovery of columns from rows and locally recoverable array codes, and for a statement like \cref{thm:T} among the
works above and those that cite them. We found neither. The search was not exhaustive; in particular we have not read
every dissertation in the area.

\subsection{Organization}\label{sec:intro-org}

\Cref{sec:prelim} fixes the model, the codes and the notation. \Cref{sec:rows,sec:tape,sec:weak} prove the reduction.
\Cref{sec:normal,sec:order,sec:entropy,sec:invariance} study the conjecture: a normal form, what order cannot do, what
entropy can, and where one half is proved. \Cref{sec:construction} is the construction side, \cref{sec:open} the open
problems, and the appendices index the Lean declarations and describe the verification.
